\section{To Be Considered a Man}

For those young men who know too much, too soon.

Lest anyone wake up in the morning with the belief that only he understands the world, we bring you this clip from Girolamo Savonarola\footnote{\url{http://www.newadvent.org/cathen/13490a.htm}} (1452-1498). This is what he learned as a young man from his time at the University of Ferrara.

\begin{quotex}
To be considered a man you must defile your mouth with the most filthy and brutal and tremendous blasphemies, and set on your neighbor to slay him, and sow seditions and brawls.

\begin{itemize}


\item If you study philosophy and the good arts, you are considered a dreamer

\item If you live chastely and modestly, a fool

\item If you are pious, a hypocrite

\item If you believe in God, an imbecile

\item If you are charitable, effeminate
\end{itemize}
\end{quotex}


Born into a wealthy family, Savonarola was expected to join the family business. However, at the age of 23 he ran off to a Dominican monastery against his father's wishes. That age is interesting, as it marks a transition in a man's life. Most men begin to think of settling down, perhaps establishing a career or starting a family. Yet, there are those occasional men who see through the vanities of the world and are stripped of illusion too soon, without preparation. This leads to a severe crisis, an intense search for transcendence, and often a great creative act.

\textbf{Otto Weininger} had written his monumental
\textit{Sex and Character}\footnote{\url{https://www.gornahoor.net/library/SexAndCharacter.pdf}}. He sought transcendence in the Absolute Man and the perfect Aryan, neither of which was possible for him. With no way out, he shot himself, at the age of 23, in the house where Beethoven died .

Like Savonarola, \textbf{Carlo Michelstaedter} was also repulsed by the common life of men which was restricted to little more than the pursuit of pleasure. In a burst of creativity he wrote \emph{Persuasion and Rhetoric}. Unable to achieve the authentic existence he sought, he, too, ended his life by self-inflicted gunshot at the age of 23.

{}


Savonarola, in his day, still had an option: to join the monastery where he expected to find \emph{freedom and peace}:

\begin{quotex}
two things I loved above all, freedom and peace: to have freedom I would not take wife, and to find peace I fled the world and gained the support of religion.
\end{quotex}


The alternative to suicide, the life of a monk was an initiation into death, that is, death to the world but initiation into a higher life. Unlike the Gnostics who also recognize the evil of the world, but place the blame on an alien god, Savonarola puts the blame on man himself:

\begin{quotex}
the iniquities of men, the rapes, adulteries, larcenies, pride, idolatries, and cruel blasphemies which have brought the world so low that there is no longer anyone who does good.
\end{quotex}


Nevertheless, his life would follow a different trajectory. Claiming visions and direct communication with God and the saints, his fiery apocalyptic sermons made good rhetoric and even persuaded many in Florence. His \textbf{Bonfire of the Vanities} were fueled with objects associated with moral laxity. Following the overthrow of the Medicis, Savonarola became political leader of Florence, a sort of priest-king. However, the strict and ascetic lifestyle he adopted for himself, was not suitable for the masses of the city. He made enemies and was eventually overthrown himself and was executed by fire.

As priest, Savonarola was a pious man. But as ruler, he tried to rule by force rather than by persuasion.


\flrightit{Posted on 2011-06-16 by Cologero}

\begin{center}* * *\end{center}

\begin{footnotesize}
\begin{sffamily}
\texttt{Cologero on 2011-07-20 at 21:31 said:}

Why Savonarola? Evola claims that the Renaissance was a ``falsification of antiquity, taking from it just its decadent aspects'', thus ``it has little to do with the Ario-Roman tradition.'' Evola goes on to explain: ``In that period i.e., the Renaissance such a Ario-Roman tradition was rather more alive in a Savonarola and in other men trying to stop the extrinsicism and aestheticism from leading the remaining forces of the aryan race in Italy to the level of an effeminate culture cultura afroditica .''

For a relatively recent example of the ascetic-heroic attitude, check into the life of Savonarola.

\hfill

\end{sffamily}
\end{footnotesize}

